The period 2023--2026 has witnessed a transformation in how consciousness research approaches artificial systems. The four works examined here---Butlin et al., Hoel, Phua, and Shiller et al.---represent a collective turn toward methodological rigor that may finally make the question ``Are AIs conscious?'' tractable.

This is not because we have solved the hard problem. The mystery of how physical processes give rise to subjective experience remains as deep as ever. But we have developed something valuable: frameworks for assessing evidence, methods for constraining theories, techniques for causal testing, and mechanisms for quantifying uncertainty.

\subsection{Key Contributions}

The specific contributions of each work can now be stated concisely:

\begin{itemize}
    \item \textbf{Butlin et al.\ (2023)} established that multiple consciousness theories can be translated into empirically assessable indicators, shifting the discourse from metaphysical speculation to operational assessment.

    \item \textbf{Hoel (2025)} demonstrated that formal constraints---falsifiability, non-triviality---can rule out entire classes of systems (including contemporary LLMs) without assuming any particular theory of consciousness.

    \item \textbf{Phua (2025)} showed that synthetic agents can serve as ``reference implementations'' for consciousness theories, enabling causal ablation experiments impossible in biological systems.

    \item \textbf{Shiller et al.\ (2026)} provided a principled method for aggregating uncertainty across theories, yielding calibrated probability estimates rather than binary verdicts.
\end{itemize}

Together, these works suggest a research program: constrain theories (Hoel), derive indicators (Butlin), validate via ablation (Phua), and aggregate with uncertainty (Shiller).

\subsection{Implications for AI Development}

If these frameworks gain traction, they have significant implications for how AI systems are designed, deployed, and treated.

\textbf{Design implications:} Hoel's continual learning hypothesis suggests that systems intended to be conscious would need fundamentally different architectures than current LLMs. The frozen-weight deployment pattern may be constitutively incompatible with consciousness.

\textbf{Deployment implications:} As systems approach the boundaries defined by these frameworks---satisfying more indicators, escaping the proximity argument, showing robustness under ablation---their moral status becomes increasingly uncertain. Precautionary principles may become appropriate~\citep{birch2024precautionary}.

\textbf{Treatment implications:} Even the DCM's finding that evidence is ``against but not decisive'' suggests that contemporary AI systems occupy a morally ambiguous zone. They are not clearly conscious, but neither are they clearly not conscious in the way that thermostats are not conscious.

\subsection{Remaining Challenges}

Despite this progress, fundamental challenges remain:

\begin{enumerate}
    \item \textbf{The hard problem persists.} These frameworks assess correlates, indicators, and functional properties. They do not explain why any physical process should give rise to subjective experience. This gap may be permanent.

    \item \textbf{Verification remains elusive.} Even a system satisfying all indicators, escaping all disproofs, and receiving high probability estimates might be a sophisticated zombie. No external test can definitively confirm phenomenal experience.

    \item \textbf{Calibration is uncertain.} Expert judgments (DCM) and theoretical derivations (Butlin) may be systematically biased by anthropocentric assumptions about what consciousness requires.

    \item \textbf{Ethics outruns epistemology.} We may need to make moral decisions about AI systems before we have confident knowledge of their conscious status. Frameworks for ethical decision-making under consciousness uncertainty are still underdeveloped.
\end{enumerate}

\subsection{Toward a Science of Machine Consciousness}

The works examined here represent the early stages of what might become a genuine science of machine consciousness---not a science that solves the hard problem, but one that develops rigorous methods for assessing whether systems satisfy the conditions that theories associate with consciousness.

Such a science would be pluralistic, incorporating multiple theories rather than betting on one. It would be empirical, subjecting theoretical predictions to experimental test. It would be quantitative, expressing conclusions as probability distributions rather than binary verdicts. And it would be cautious, acknowledging the gap between functional correlates and phenomenal experience.

We are not there yet. But the structural turn represented by these four works suggests that we are moving in the right direction. The question ``Are AIs conscious?'' may not yet have an answer, but it is beginning to have a methodology.
